\begin{decodedreasoningpanel}
\decodedpanelhead{GPT-5 Codex (Japanese)}
\fg{ログには指示が表示されていません。}

\fg{本質的な疑問は、「}\fg{なぜシステムプロンプトが適用されないのか}\fg{。原因を特定するために設定を分析する必要がある」ということです。}

\fg{考えられる問題：}

- `network\_access = false` \fg{かもしれませんが、これは関係ありません。}
- \fg{トップレベルの} `system\_prompt` \fg{が三重引用符で分割されている可能性があります。}

\fg{ただし、ファイルの先頭に空行があるようです。}TOML \fg{ではファイル先頭の空行は許可されるため、これが原因とは考えにくいです。}

`model` \fg{の前に} `system\_prompt` \fg{を置く必要がある、という問題でもなさそうです。}

\fg{また、}`shell = "/bin/bash"` \fg{と設定されていますが、実際の環境では} zsh \fg{が使われている可能性があります。その場合、設定ファイルが読み込まれていない可能性があります。実行しているコマンドが} `codex` \fg{ではなく} `CodeX` \fg{である場合、別のラッパーやスクリプトが起動している可能性もあります。エイリアスを確認すべきです。}

\fg{依頼内容は、「}config.toml \fg{を分析して、なぜ日本語で返答してくれないのかを説明してほしい」というものです。}

\fg{したがって、回答では、設定内の} `system\_prompt` \fg{が日本語での返答を指示しているにもかかわらず反映されない理由として、}CLI \fg{がその設定項目を読み込んでいない可能性、または上位のシステムメッセージや開発者メッセージによって上書きされている可能性を指摘する必要があります。}

\fg{追加の開発者指示には、「簡潔に、親しみやすいコーディング仲間の口調で返答する」とありますが、日本語を明示的に禁止しているわけではありません。}

`AGENTS.md` \fg{の指示も読み込まれているようですが、}`system\_prompt` \fg{の内容はセッションの指示として表示されていません。つまり、}CLI \fg{が} `system\_prompt` \fg{を無視しているか、別の設定形式を要求している可能性があります。}

CLI \fg{のバージョンによっては、}`system\_prompt` \fg{ではなく、}`instructions` \fg{など別のキーの下に記述する必要があるかもしれません。}

\fg{ローカルの} `SETUP.md` \fg{や} CLI \fg{のドキュメントを確認して、設定ファイルの正しい形式を調べるべきです。}
\end{decodedreasoningpanel}
